\documentclass[11pt,a4paper]{article}
% preamble.tex — estilo común de todos los apuntes Froth.
% Cada apunte hace % preamble.tex — estilo común de todos los apuntes Froth.
% Cada apunte hace % preamble.tex — estilo común de todos los apuntes Froth.
% Cada apunte hace \input{preamble} justo después de \documentclass.
\usepackage[margin=2.4cm]{geometry}
\usepackage{xcolor}
\definecolor{litblue}{RGB}{31,79,121}
\definecolor{litgreen}{RGB}{27,94,32}
\definecolor{litorange}{RGB}{230,81,0}
\usepackage{parskip}
\usepackage{titlesec}
\titleformat{\section}{\Large\bfseries\color{litblue}}{\thesection}{0.6em}{}
\titleformat{\subsection}{\large\bfseries\color{litblue}}{\thesubsection}{0.6em}{}
\usepackage{tcolorbox}
\usepackage{fancyhdr}
\pagestyle{fancy}
\fancyhf{}
\fancyhead[L]{\small\color{litblue}Froth — Apuntes de ML}
\fancyhead[R]{\small\thepage}
\renewcommand{\headrulewidth}{0.4pt}
\usepackage[colorlinks=true,linkcolor=litblue,urlcolor=litblue]{hyperref}

% Cabecera de cada apunte: \cabecera{numero}{titulo}
\newcommand{\cabecera}[2]{%
  \begin{tcolorbox}[colback=litblue,coltext=white,colframe=litblue,arc=2mm]
    {\Large\bfseries Apunte #1 — #2}\\[3pt]
    {\small Proyecto Froth · aprender Machine Learning construyendo}
  \end{tcolorbox}\vspace{0.8em}}

% Cajas temáticas
\newtcolorbox{concepto}[1]{colback=blue!4,colframe=litblue,title={Concepto: #1},arc=1.5mm}
\newtcolorbox{practica}{colback=green!5,colframe=litgreen,title={Pruébalo tú (teclado en mano)},arc=1.5mm}
\newtcolorbox{ojo}{colback=orange!8,colframe=litorange,title={Ojo / error típico},arc=1.5mm}
\newtcolorbox{resumen}{colback=gray!8,colframe=gray!60,title={En una frase},arc=1.5mm}

\newcommand{\codigo}[1]{\texttt{#1}}
 justo después de \documentclass.
\usepackage[margin=2.4cm]{geometry}
\usepackage{xcolor}
\definecolor{litblue}{RGB}{31,79,121}
\definecolor{litgreen}{RGB}{27,94,32}
\definecolor{litorange}{RGB}{230,81,0}
\usepackage{parskip}
\usepackage{titlesec}
\titleformat{\section}{\Large\bfseries\color{litblue}}{\thesection}{0.6em}{}
\titleformat{\subsection}{\large\bfseries\color{litblue}}{\thesubsection}{0.6em}{}
\usepackage{tcolorbox}
\usepackage{fancyhdr}
\pagestyle{fancy}
\fancyhf{}
\fancyhead[L]{\small\color{litblue}Froth — Apuntes de ML}
\fancyhead[R]{\small\thepage}
\renewcommand{\headrulewidth}{0.4pt}
\usepackage[colorlinks=true,linkcolor=litblue,urlcolor=litblue]{hyperref}

% Cabecera de cada apunte: \cabecera{numero}{titulo}
\newcommand{\cabecera}[2]{%
  \begin{tcolorbox}[colback=litblue,coltext=white,colframe=litblue,arc=2mm]
    {\Large\bfseries Apunte #1 — #2}\\[3pt]
    {\small Proyecto Froth · aprender Machine Learning construyendo}
  \end{tcolorbox}\vspace{0.8em}}

% Cajas temáticas
\newtcolorbox{concepto}[1]{colback=blue!4,colframe=litblue,title={Concepto: #1},arc=1.5mm}
\newtcolorbox{practica}{colback=green!5,colframe=litgreen,title={Pruébalo tú (teclado en mano)},arc=1.5mm}
\newtcolorbox{ojo}{colback=orange!8,colframe=litorange,title={Ojo / error típico},arc=1.5mm}
\newtcolorbox{resumen}{colback=gray!8,colframe=gray!60,title={En una frase},arc=1.5mm}

\newcommand{\codigo}[1]{\texttt{#1}}
 justo después de \documentclass.
\usepackage[margin=2.4cm]{geometry}
\usepackage{xcolor}
\definecolor{litblue}{RGB}{31,79,121}
\definecolor{litgreen}{RGB}{27,94,32}
\definecolor{litorange}{RGB}{230,81,0}
\usepackage{parskip}
\usepackage{titlesec}
\titleformat{\section}{\Large\bfseries\color{litblue}}{\thesection}{0.6em}{}
\titleformat{\subsection}{\large\bfseries\color{litblue}}{\thesubsection}{0.6em}{}
\usepackage{tcolorbox}
\usepackage{fancyhdr}
\pagestyle{fancy}
\fancyhf{}
\fancyhead[L]{\small\color{litblue}Froth — Apuntes de ML}
\fancyhead[R]{\small\thepage}
\renewcommand{\headrulewidth}{0.4pt}
\usepackage[colorlinks=true,linkcolor=litblue,urlcolor=litblue]{hyperref}

% Cabecera de cada apunte: \cabecera{numero}{titulo}
\newcommand{\cabecera}[2]{%
  \begin{tcolorbox}[colback=litblue,coltext=white,colframe=litblue,arc=2mm]
    {\Large\bfseries Apunte #1 — #2}\\[3pt]
    {\small Proyecto Froth · aprender Machine Learning construyendo}
  \end{tcolorbox}\vspace{0.8em}}

% Cajas temáticas
\newtcolorbox{concepto}[1]{colback=blue!4,colframe=litblue,title={Concepto: #1},arc=1.5mm}
\newtcolorbox{practica}{colback=green!5,colframe=litgreen,title={Pruébalo tú (teclado en mano)},arc=1.5mm}
\newtcolorbox{ojo}{colback=orange!8,colframe=litorange,title={Ojo / error típico},arc=1.5mm}
\newtcolorbox{resumen}{colback=gray!8,colframe=gray!60,title={En una frase},arc=1.5mm}

\newcommand{\codigo}[1]{\texttt{#1}}

\begin{document}
\cabecera{28}{Conectores multi-fuente: tu cuenta institucional entra al juego}

\section{Por qué no basta OpenAlex}

OpenAlex es enorme ($\sim$250M trabajos) pero no lo tiene todo, y a veces le faltan
abstracts. La idea de Batman: que el usuario \textbf{conecte sus propias cuentas} (estilo
``Connect with Google'') y vea papers que la base gratuita no trae.

\section{Los conectores (froth/sources.py)}

\begin{center}
\begin{tabular}{lll}
\textbf{Fuente} & \textbf{Requiere} & \textbf{Qué aporta} \\
\hline
OpenAlex & nada (key opcional) & la BASE: abstracts + referencias \\
Crossref & solo un email & metadatos masivos de editoriales \\
Semantic Scholar & key GRATIS (opcional) & abstracts + PDFs open access \\
Scopus (Elsevier) & key institucional & el catálogo curado de Elsevier \\
ResearchGate & \multicolumn{2}{l}{\textbf{NO}: sin API pública; su ToS prohíbe scraping} \\
\end{tabular}
\end{center}

Principios de diseño:
\begin{itemize}
  \item \textbf{OpenAlex manda}: es el registro preferido (el más rico); los extras solo
        AÑADEN papers que no teníamos.
  \item \textbf{Dedup por DOI y luego por título} normalizado — el mismo paper llega por
        varias fuentes y debe contar una vez.
  \item \textbf{Fallo suave}: si una fuente está caída o te limita, avisa y el pull sigue
        con las demás. Nunca un conector tumba la cosecha.
  \item \textbf{Las keys son tuyas y locales}: viven en \codigo{.froth\_keys} (gitignored).
        En la web pública vivirán solo en tu sesión de navegador — Froth jamás custodia
        keys ajenas en un servidor.
\end{itemize}

\section{La primera cosecha multi-fuente (números reales)}

\begin{itemize}
  \item Crossref: 227 resultados crudos $\rightarrow$ el dedup reconoció \textbf{58
        duplicados} de OpenAlex y aportó \textbf{169 nuevos} (tras el filtro de abstract:
        corpus 126 $\rightarrow$ \textbf{157}).
  \item Semantic Scholar sin key: rate-limit inmediato $\rightarrow$ \textbf{falló suave} tal
        como se diseñó (con la key gratuita queda estable).
  \item El corpus ampliado reveló un \textbf{subtema emergente}: física de
        burbujas/partículas finas (+ bastnäsita) — antes enterrado, ahora cluster propio.
        Más datos $\rightarrow$ más resolución del mapa.
  \item El filtro de relevancia solo descartó 6: la cosecha extra venía limpia — el
        auto-lavado del Apunte 17 escaló sin tocarlo.
\end{itemize}

\begin{resumen}
sources.py suma Crossref/S2/Scopus a OpenAlex con dedup DOI→título y fallo suave; las keys
son locales y del usuario (panel ``Connect institutional account''). Primera cosecha: +31
papers netos y un subtema nuevo emergió. ResearchGate queda fuera por diseño (sin API legal).
\end{resumen}

\end{document}
